%% LyX 2.5.0 created this file.  For more info, see https://www.lyx.org/.
%% Do not edit unless you really know what you are doing.
\documentclass[12pt,croatian]{article}
\usepackage[T1]{fontenc}
\usepackage[latin2]{inputenc}
\usepackage{setspace}
\onehalfspacing

\makeatletter
%%%%%%%%%%%%%%%%%%%%%%%%%%%%%% User specified LaTeX commands.
% Added by lyx2lyx
\setlength{\parskip}{\medskipamount}
\setlength{\parindent}{0pt}

\makeatother

\usepackage{babel}
\begin{document}

\section*{Odre�ivanje pada brzine broda na valovitom moru}

\subsection*{Pojmovi}
\begin{itemize}
\item $R_{CW}$ - ukupni otpor broda na mirnoj vodi (uklju�uje otpor trupa,
privjesci, zraka, ...)
\item $R_{Wind}$ - dodatni otpor vjetra
\item $\overline{R_{W}}$ - osrednjeni dodatni otpor uslijed valova
\item $P_{B}$ - snaga motora na ko�nici
\item $V$ - brzina broda na mirnom moru
\item $V_{W}$ - brzina broda na valovitom moru
\item $\Delta V$ - pad brzine
\end{itemize}

\subsection*{Metodologija}

Za brod u plovidbi na odre�enom stanju mora, ukupni otpor zbroj je
svih komponenti otpora:
\[
R_{T}=R_{CW}+R_{Wind}+\overline{R_{W}}.
\]
Pad brzine mo�e se izra�unati ukoliko se pretpostavi da se koristi
jednaka snaga na mirnoj vodi kao i pri plovidbi na odre�enom stanju
valovitog mora:
\[
P_{B}\left(V\right)=P_{B,W}\left(V_{W}\right),
\]
jer je poznata snaga $P_{B}$ za mirno more i projektnu brzinu $V$.
Pad brzine je definiran kao:
\[
\Delta V=V-V_{W}.
\]
Najbr�a aproksimacija mo�e se dobiti ako pretpostavimo da je iskoristivost
propulzije jednaka na mirnoj vodi odnosu na valovito more (iako nije).
Tada mo�emo izjedna�iti ne samo snage motora (kao u jednad�bi gore),
ve� i efektivne snage $P_{E}$, �to je produkt otpora i brzine:
\[
R_{CW}\left(V\right)\,V=R_{T}\left(V_{W}\right)\,V_{W}.
\]
Problem je �to je brzina $V_{W}$ trenutno nepoznata. Dolje je dana
procedura prora�una pada brine bez kori�tenja derivacije krivulja
otpora.

\subsection*{Procedura}
\begin{enumerate}
\item Izra�unati krivulju otpora na mirnoj vodi $R_{CW}$, npr. pomo�u Maxsurf
Resistance + privjesci i sl.
\item Izra�unati krivulju dodatnog otpora na odre�enom stanju mora $\overline{R_{W}}$,
npr. prema Maxsurf Motions
\item Izra�unati krivulju otpora vjetra na odre�enom stanju mora $R_{Wind}$,
prema brzini vjetra za zadano stanje, projiciranoj povr�ini broda
i koeficijentu otpora.
\item Zbrojiti krivulje u krivulju ukupnog otpora na odre�enom stanju mora,
$R_{T}$
\item Izra�unati brzinu uvo�enjem aproksimacije $V\approx V_{W}$, tj. $\Delta V\rightarrow0$,
o�itavanjem otpora $R_{CW}$ i $R_{T}$ na brzini $V$: 
\[
V_{W0}=V\,\frac{R_{CW}\left(V\right)}{R_{T}\left(V\right)}
\]
\item Za dobivenu brzinu ponoviti prora�un, o�itavanjem otpora $R_{T}$
na gore izra�unatoj brzini $V_{W0}$:
\[
V_{W}=V\,\frac{R_{CW}\left(V\right)}{R_{T}\left(V_{W0}\right)}
\]
\item Korak 6 mo�ete ponoviti vi�e puta za bolju konvergenciju.
\item Izra�unati pad brzine: $\Delta V=V-V_{W}$.
\end{enumerate}

\subsection*{Zadatak}

Prema danoj proceduri i zadanom brodu potrebno je izra�unati pad brzine
broda uslijed valova, ukoliko brod ina�e plovi brzinom od 17 �vorova
na mirnom moru. Razmatrati valove u pramac i jednoparametarski spektar
za koji zna�ajna valna visina iznosi 3 m.
\end{document}
